\documentclass[12pt,a4paper,headings=optiontohead]{scrbook}
\usepackage[utf8]{inputenc}
\usepackage[italian]{babel}
\usepackage[margin=1.8cm,bottom=7em]{geometry}
\usepackage[subpreambles=false]{standalone}
\usepackage{amsmath}
\usepackage{amssymb}
\usepackage{amsthm} 
\usepackage{cancel}
\usepackage{graphicx}
\usepackage{mathtools}
\usepackage{float}
\usepackage{enumitem}
\usepackage{algorithm}
\usepackage{algorithmic}


\usepackage[normalem]{ulem}

\usepackage{blkarray}% http://ctan.org/pkg/blkarray

\newcommand{\matindex}[1]{\mbox{\scriptsize#1}}% Matrix index

\renewcommand{\qedsymbol}{\rule{0.7em}{0.7em}}
\DeclarePairedDelimiter{\abs}{\lvert}{\rvert}
\newcommand{\inter}{\begin{matrix}\prod\end{matrix}}
\newcommand{\verteq}{\rotatebox{90}{$\,=$}}
\newcommand{\equalto}[2]{\underset{\scriptstyle\overset{\mkern4mu\verteq}{#2}}{#1}}
\DeclarePairedDelimiter{\norma}{\lVert}{\rVert}
\newtheorem*{esempio}{Esempio}

\usepackage{import}
\usepackage{hyperref}
\begin{document}

%------------------------------------------------------------------------------------------------------
%------------------------------------------------------------------------------------------------------
%-----------------------------------------------INTESTAZIONE-------------------------------------------
%------------------------------------------------------------------------------------------------------
%------------------------------------------------------------------------------------------------------

\begin{titlepage}

\newcommand{\HRule}{\rule{\linewidth}{0.5mm}} % Defines a new command for the horizontal lines, change thickness here

\center % Center everything on the page
 
%----------------------------------------------------------------------------------------
%	HEADING SECTIONS
%----------------------------------------------------------------------------------------


\large Dalle dispense del\\[0.5cm] % Major heading such as course name
\textsc{\Large Prof. Marco Vianello}\\[1.5cm] % Minor heading such as course title

\includegraphics[scale=0.13]{maths.png}\\[0.3cm]
%----------------------------------------------------------------------------------------
%	TITLE SECTION
%----------------------------------------------------------------------------------------

\HRule \\[0.4cm]
{ \huge \bfseries Appunti di}\\
{ \huge \bfseries Calcolo Numerico\\[0.15 cm]} % Title of your document
\HRule \\[1.5cm]
 
%----------------------------------------------------------------------------------------
%	AUTHOR SECTION
%----------------------------------------------------------------------------------------

\begin{minipage}{0.4\textwidth}
\begin{flushleft} \large
\emph{\Large{Autori:}}\\
\quad \quad Michele \textsc{Veronesi}\\
\quad \quad Silvia \textsc{Bazzeato}\\
\quad \quad Simone \textsc{De Renzis}\\
\quad \quad Michele \textsc{Baldisseri}\\
\quad \quad Andrea \textsc{Cecchin}\\
\quad \quad Stefano \textsc{Rizzo}\\
\quad \quad Leonardo \textsc{Tredese}\\
\quad \quad Alessandro \textsc{Pirolo}\\
\quad \quad Filippo \textsc{Pinton}
\end{flushleft}

\end{minipage}\\[3cm]

%----------------------------------------------------------------------------------------
%	DATE SECTION
%----------------------------------------------------------------------------------------

\LARGE Università degli Studi di Padova\\[0.4cm] % Name of your university/college
\textsc{\large Dipartimento di Matematica}\\[0.05cm]
\textsc{\large Corso di Laurea in Informatica}\\[1cm] % Include a department/university logo - this will require the graphicx package

{\Large Anno accademico 2019 - 2020}\\[2cm] % Date, change the \today to a set date if you want to be precise

\vfill % Fill the rest of the page with whitespace

\end{titlepage}


\begin{center}
\pagebreak

\section*{Premessa}
\begin{minipage}{0.9\textwidth} \large

Dagli appunti raccolti durante un anno costellato di difficoltà, è nata l'esigenza di unificare il dettagliato materiale in un documento quanto più possibile uniforme e fruibile. Chi dedicatosi a una cosa e chi ad un'altra, la trascrizione è frutto della collaborazione di studenti: non ci si aspetti una forma perfetta e si studi (come sempre del resto) con occhio critico e attento agli errori.

\end{minipage}

\end{center}
\pagebreak

\tableofcontents

% ############ CAPITOLO 1 ##################
\chapter{Sistema floating-point e propagazione degli errori; costo computazionale}
\chaptermark{Floating-point, errori, costo computazionale}

\import{lezione1/}{lezione1.tex}
\newpage

\import{lezione2/}{lezione2.tex}
\newpage

\import{lezione3/}{lezione3.tex}
\newpage

\import{lezione4/}{lezione4.tex}
\newpage

\import{lezione5/}{lezione5.tex}
\newpage

\import{lezione6/}{lezione6.tex}
\newpage

\import{lezione7/}{lezione7.tex}
\newpage

% ############# CAPITOLO 2 #####################
\chapter{Soluzione numerica di equazioni non lineari}

\import{lezione8/}{lezione8.tex}
\newpage

\import{lezione9/}{lezione9.tex}
\newpage

\import{lezione10/}{lezione10.tex}
\newpage

\import{lezione11/}{lezione11.tex}
\newpage

% ############ CAPITOLO 3 ######################
\chapter{Interpolazione e approssimazione di dati e funzioni}

\import{lezione12/}{lezione12.tex}
\newpage

\import{lezione13/}{lezione13.tex}
\newpage

\import{lezione14/}{lezione14.tex}
\newpage

\import{lezione15/}{lezione15.tex}
\newpage

% ############## CAPITOLO 4 #####################
\chapter{Integrazione numerica e derivazione numerica}

\import{lezione16/}{lezione16.tex}
\newpage

\import{lezione17/}{lezione17.tex}
\newpage

\import{lezione18/}{lezione18.tex}
\newpage

% ################ CAPITOLO 5 #####################
\chapter{Elementi di algebra lineare numerica}

\import{lezione19/}{lezione19.tex}
\newpage

\import{lezione20/}{lezione20.tex}
\newpage

\import{lezione21/}{lezione21.tex}
\newpage

\import{lezione22/}{lezione22.tex}
\newpage

\import{lezione23/}{lezione23.tex}
\newpage

\import{lezione24/}{lezione24.tex}
\newpage
\end{document}
